\documentclass[11pt]{article}

\usepackage[margin=1in]{geometry}
\usepackage{amsmath,amssymb,amsthm}
\usepackage{booktabs}
\usepackage{longtable}
\usepackage{array}
\usepackage{hyperref}
\usepackage{listings}
\usepackage{xcolor}

\hypersetup{
    colorlinks=true,
    linkcolor=blue,
    citecolor=blue,
    urlcolor=blue
}

\lstset{
    basicstyle=\ttfamily\small,
    breaklines=true,
    frame=single,
    columns=fullflexible,
    keepspaces=true
}

\newcommand{\pig}{\pi_g}
\newcommand{\pigN}{\pi_g^{(N)}}
\newcommand{\QN}{Q_N}

\newtheorem{theorem}{Theorem}
\newtheorem{proposition}{Proposition}
\newtheorem{remark}{Remark}

\title{Classical Appearances of \texorpdfstring{\(\pi_g\)}{pi-g}:\\
A Geometric Substitution Note}

\author{Wayne Baker}

\date{July 5, 2026}

\begin{document}

\maketitle

\begin{abstract}
This short note records representative classical formulae in which the usual
constant \(\pi\) is rewritten as a geometric limiting constant \(\pi_g\). The
point is not to claim new classical identities. Rather, once a prior geometric
argument establishes
\[
\pi_g=\pi,
\]
every exact occurrence of \(\pi\) may be equivalently written as \(\pi_g\).
The distinction between the completed limiting value \(\pi_g\) and finite
approximants \(\pi_g^{(N)}\) is emphasized. A numerical consistency test is
included using an \(N\)-series polygonal approximation to \(\pi_g\) and the
Gaussian integral.
\end{abstract}

\section{Purpose of the note}

Suppose a geometric construction or polygonal acceleration scheme defines a
sequence of finite approximants
\[
\pi_g^{(N)}
\]
and proves that these approximants converge to the ordinary circle constant:
\[
\pi_g
=
\lim_{N\to\infty}\pi_g^{(N)}
=
\pi.
\]

Then \(\pi_g\) is not a new numerical value different from \(\pi\). It is an
exact geometric representation of the same completed constant. Consequently,
any exact classical identity involving \(\pi\) may be rewritten with
\(\pi_g\).

The main distinction is therefore:

\[
\pi_g^{(N)} \text{ is finite and approximate,}
\]

while

\[
\pi_g \text{ is exact, with } \pi_g=\pi.
\]

Thus exact substitutions use \(\pi_g\), whereas numerical tests necessarily use
finite approximants \(\pi_g^{(N)}\).

\section{The substitution principle}

\begin{proposition}[Geometric substitution principle]
Assume that a geometric limiting constant \(\pi_g\) has been established by
\[
\pi_g=\pi.
\]
Then every exact identity of the form
\[
F(\pi)=A
\]
may be rewritten equivalently as
\[
F(\pi_g)=A,
\]
provided the expression \(F\) is evaluated on the same branch and domain as in
the original identity.
\end{proposition}

\begin{proof}
Since \(\pi_g=\pi\), substitution gives
\[
F(\pi_g)=F(\pi).
\]
Therefore, if \(F(\pi)=A\), then \(F(\pi_g)=A\). This is equality by
replacement of identical constants, not a new analytic theorem.
\end{proof}

\begin{remark}
The finite approximation \(\pi_g^{(N)}\) should not be used in exact identities
unless an error term is included. For finite \(N\),
\[
\pi_g^{(N)}\approx \pi_g=\pi,
\]
but generally
\[
\pi_g^{(N)}\neq \pi.
\]
Therefore
\[
F(\pi_g^{(N)})\approx F(\pi),
\]
whereas
\[
F(\pi_g)=F(\pi)
\]
is exact.
\end{remark}

\section{Representative classical formulae}

The following table records standard appearances of \(\pi\) and their
equivalent \(\pi_g\)-forms. These entries illustrate substitution sites, not
independent numerical verifications; several (e.g.\ the Gamma value) follow
directly from the Gaussian integral entry via standard identities rather than
constituting separate checks. The single worked numerical example is given in
Section~\ref{sec:gaussian-test}.


\begin{center}
\small
\renewcommand{\arraystretch}{1.45}
\begin{longtable}{>{\raggedright\arraybackslash}p{0.20\textwidth}
                  >{\centering\arraybackslash}p{0.35\textwidth}
                  >{\centering\arraybackslash}p{0.35\textwidth}}
\toprule
\textbf{Identity} & \textbf{Classical form} & \textbf{\(\pi_g\)-form} \\
\midrule
\endfirsthead
\toprule
\textbf{Identity} & \textbf{Classical form} & \textbf{\(\pi_g\)-form} \\
\midrule
\endhead
\bottomrule
\endfoot

Euler identity
& \(\displaystyle e^{i\pi}+1=0\)
& \(\displaystyle e^{i\pi_g}+1=0\) \\

Gaussian integral
& \(\displaystyle \int_{-\infty}^{\infty}e^{-x^2}\,dx=\sqrt{\pi}\)
& \(\displaystyle \int_{-\infty}^{\infty}e^{-x^2}\,dx=\sqrt{\pi_g}\) \\

Basel problem
& \(\displaystyle \sum_{n=1}^{\infty}\frac{1}{n^2}=\frac{\pi^2}{6}\)
& \(\displaystyle \sum_{n=1}^{\infty}\frac{1}{n^2}=\frac{\pi_g^2}{6}\) \\

Gamma value
& \(\displaystyle \Gamma\!\left(\frac12\right)=\sqrt{\pi}\)
& \(\displaystyle \Gamma\!\left(\frac12\right)=\sqrt{\pi_g}\) \\

Euler reflection formula
& \(\displaystyle \Gamma(z)\Gamma(1-z)=\frac{\pi}{\sin(\pi z)}\)
& \(\displaystyle \Gamma(z)\Gamma(1-z)=\frac{\pi_g}{\sin(\pi_g z)}\) \\

Wallis product
& \(\displaystyle \frac{\pi}{2}=\prod_{n=1}^{\infty}\frac{(2n)(2n)}{(2n-1)(2n+1)}\)
& \(\displaystyle \frac{\pi_g}{2}=\prod_{n=1}^{\infty}\frac{(2n)(2n)}{(2n-1)(2n+1)}\) \\

Leibniz series
& \(\displaystyle \frac{\pi}{4}=1-\frac13+\frac15-\frac17+\cdots\)
& \(\displaystyle \frac{\pi_g}{4}=1-\frac13+\frac15-\frac17+\cdots\) \\

Fourier-Gaussian identity
& \(\displaystyle \begin{gathered}\int_{-\infty}^{\infty} e^{-\pi x^2}e^{-2\pi i x\xi}\,dx\\=e^{-\pi \xi^2}\end{gathered}\)
& \(\displaystyle \begin{gathered}\int_{-\infty}^{\infty} e^{-\pi_g x^2}e^{-2\pi_g i x\xi}\,dx\\=e^{-\pi_g \xi^2}\end{gathered}\) \\

Stirling formula
& \(\displaystyle n!\sim \sqrt{2\pi n}\left(\frac ne\right)^n\)
& \(\displaystyle n!\sim \sqrt{2\pi_g n}\left(\frac ne\right)^n\) \\

Cauchy integral formula
& \(\displaystyle f(a)=\frac{1}{2\pi i}\oint_{\gamma}\frac{f(z)}{z-a}\,dz\)
& \(\displaystyle f(a)=\frac{1}{2\pi_g i}\oint_{\gamma}\frac{f(z)}{z-a}\,dz\) \\
\end{longtable}
\end{center}

Each formula in the right column is exactly equivalent to its classical
counterpart once
\[
\pi_g=\pi
\]
has been established.

\section{Finite approximants and why small numerical differences appear}

The exact identity is
\[
\pi_g=\pi.
\]

However, a computer does not evaluate the completed infinite limit directly.
Instead it evaluates a finite approximation
\[
\pi_g^{(N)}.
\]

A typical polygonal quantity is
\[
Q_N=\frac{Nc_N}{2},
\]
where \(c_N\) is the chord length of a regular \(N\)-gon in the unit circle.
The first \(N\)-series correction is
\[
\widehat Q_N
=
\frac{4Q_{2N}-Q_N}{3}.
\]

Higher correction levels may be obtained by repeated even-power cancellation.
For example, if
\[
R_N^{(0)}=Q_N,
\]
then define
\[
R_N^{(m)}
=
\frac{4^m R_{2N}^{(m-1)}-R_N^{(m-1)}}{4^m-1}.
\]

A finite-stage geometric approximation to \(\pi_g\) is then
\[
\pi_g^{(N)}=R_N^{(m)}
\]
for some finite \(N\) and correction order \(m\).

Thus the exact equality is
\[
\pi_g
=
\lim_{N\to\infty}\pi_g^{(N)}
=
\pi,
\]
while the numerical computation only observes
\[
\pi_g^{(N)}\approx \pi.
\]

Small differences in numerical tables arise from three sources:

\[
\text{finite }N,
\]

\[
\text{finite floating-point precision,}
\]

and, for integral tests,

\[
\text{finite integration windows.}
\]

For example, a numerical Gaussian integral may use
\[
\int_{-6}^{6}e^{-x^2}\,dx
\]
instead of the exact improper integral
\[
\int_{-\infty}^{\infty}e^{-x^2}\,dx.
\]

The missing Gaussian tail is very small, but it is not literally zero. Concretely,
\[
2\int_{6}^{\infty}e^{-x^2}\,dx = \sqrt{\pi}\,\operatorname{erfc}(6)
\approx 3.81\times10^{-17},
\]
which is roughly two orders of magnitude smaller than the
\(3\times10^{-15}\) discrepancy observed in the normalized-area test below.
The truncation error is therefore not the dominant contributor to that
discrepancy; ordinary double-precision floating-point roundoff (accumulated
over the chord-doubling recursion, the N-series acceleration, and the Simpson
quadrature) accounts for essentially all of it.

\section{Gaussian normalization test}
\label{sec:gaussian-test}

The classical Gaussian integral is
\[
\int_{-\infty}^{\infty}e^{-x^2}\,dx=\sqrt{\pi}.
\]

Using
\[
\pi_g=\pi,
\]
this becomes
\[
\int_{-\infty}^{\infty}e^{-x^2}\,dx=\sqrt{\pi_g}.
\]

Equivalently, the normalized geometric Gaussian is
\[
G_g(x)=\frac{e^{-x^2}}{\sqrt{\pi_g}},
\]
and its exact area is
\[
\int_{-\infty}^{\infty}G_g(x)\,dx=1.
\]

A finite numerical test replaces \(\pi_g\) by \(\pi_g^{(N)}\) and replaces
\((-\infty,\infty)\) by a large finite interval:
\[
\frac{1}{\sqrt{\pi_g^{(N)}}}
\int_{-6}^{6}e^{-x^2}\,dx
\approx 1.
\]

Using a second-order \(N\)-series value computed from
\[
N=24576,\quad 49152,\quad 98304,
\]
one obtains the following sample values.

\begin{center}
\begin{tabular}{lll}
\toprule
\textbf{Quantity} & \textbf{Computed value} & \textbf{Reference/comparison} \\
\midrule
\(Q_{24576}\) & \(3.14159264503368973\) & -- \\
\(Q_{49152}\) & \(3.14159265145076638\) & -- \\
\(Q_{98304}\) & \(3.14159265305503599\) & -- \\
\(\pi_g^{(N)}\) & \(3.14159265358979267\) & \(\pi\approx 3.14159265358979312\) \\
\(\sqrt{\pi_g^{(N)}}\) & \(1.77245385090551588\) & \(\sqrt{\pi}\approx 1.77245385090551588\) \\
\(\int_{-6}^{6}e^{-x^2}\,dx\) & \(1.77245385090551033\) & \(\sqrt{\pi_g^{(N)}}\) \\
Difference & \(-5.551115123125783\times 10^{-15}\) & \(0\) in the exact limit \\
Normalized area & \(0.99999999999999689\) & \(1\) in the exact limit \\
\bottomrule
\end{tabular}
\end{center}

The computed normalized area is
\[
0.99999999999999689,
\]
which differs from \(1\) by approximately
\[
3.11\times 10^{-15}.
\]

This is consistent with ordinary double-precision numerical error, as noted
above; the finite integration window contributes negligibly by comparison.

\section{Interpretation}

The numerical experiment should be interpreted carefully.

It does not prove a new form of the Gaussian integral. The classical theorem
already states
\[
\int_{-\infty}^{\infty}e^{-x^2}\,dx=\sqrt{\pi}.
\]

Because
\[
\pi_g=\pi,
\]
the corresponding \(\pi_g\)-form is the same theorem written through a
geometric limiting constant:
\[
\int_{-\infty}^{\infty}e^{-x^2}\,dx=\sqrt{\pi_g}.
\]

The numerical calculation simply verifies that a high-accuracy finite
geometric approximation \(\pi_g^{(N)}\) behaves as expected when substituted
into a classical formula involving \(\pi\).

Thus the purpose of the note is representational:

\[
\text{exact appearances of } \pi
\text{ may be rewritten as exact appearances of } \pi_g.
\]

The mathematics is unchanged, but the representation emphasizes that the same
constant may be obtained as a completed geometric limit.

\section{Conclusion}

Once
\[
\pi_g=\pi
\]
has been established, the substitutions in this note are exact. The finite
values \(\pi_g^{(N)}\) used in numerical tests are approximations to the
completed constant, not replacements for it in exact identities.

Therefore the proper hierarchy is
\[
\pi_g^{(N)}
\approx
\pi_g
=
\pi.
\]

The formulae listed above show that the geometric constant \(\pi_g\), once
identified with \(\pi\), travels consistently through classical analysis,
number theory, special functions, asymptotics, Fourier analysis, and complex
analysis.

\appendix

\section{Python verification script}

The following script computes \(\pi_g^{(N)}\) from the polygonal/N-series
expression and uses it in a Gaussian normalization test.

\begin{lstlisting}[language=Python]
import math


def next_chord_by_doubling(c):
    """
    Given c_N = chord length of a regular N-gon in the unit circle,
    compute c_{2N} using only square roots.

    Stable form:
        c_{2N} = c_N / sqrt(2 + 2 sqrt(1 - c_N^2/4))
    """
    return c / math.sqrt(2.0 + 2.0 * math.sqrt(1.0 - 0.25 * c * c))


def geometric_Q_values(start_doublings, extra_values):
    """
    Start from the regular hexagon:
        N = 6, c_6 = 1.

    After start_doublings, return Q_N, Q_{2N}, Q_{4N}, ...
    where
        Q_N = N c_N / 2.
    """
    N = 6
    c = 1.0

    for _ in range(start_doublings):
        c = next_chord_by_doubling(c)
        N *= 2

    values = []

    for _ in range(extra_values):
        Q = 0.5 * N * c
        values.append((N, Q))

        c = next_chord_by_doubling(c)
        N *= 2

    return values


def n_series_acceleration(Q_values, order):
    """
    Even-power Richardson/N-series acceleration.

    order = 0:
        Q_N

    order = 1:
        (4 Q_{2N} - Q_N) / 3

    order = 2:
        applies the next correction level as well.
    """
    current = list(Q_values)

    for r in range(1, order + 1):
        factor = 4.0 ** r
        current = [
            (factor * current[i + 1] - current[i]) / (factor - 1.0)
            for i in range(len(current) - 1)
        ]

    return current[0]


def compute_pi_g(start_doublings=12, order=2):
    """
    Compute pi_g from the geometric polygon/N-series expression.

    start_doublings=12 gives N = 6 * 2^12 = 24576.
    order=2 uses Q_N, Q_{2N}, Q_{4N}.
    """
    values = geometric_Q_values(
        start_doublings=start_doublings,
        extra_values=order + 1
    )

    Q_values = [Q for _, Q in values]
    pi_g = n_series_acceleration(Q_values, order=order)

    return pi_g, values


def simpson(f, a, b, n):
    """
    64-bit Simpson integration.
    """
    if n % 2:
        n += 1

    h = (b - a) / n
    s = f(a) + f(b)

    for i in range(1, n):
        x = a + i * h
        s += (4.0 if i % 2 else 2.0) * f(x)

    return s * h / 3.0


pi_g, source_values = compute_pi_g(start_doublings=12, order=2)

L = 6.0
n_steps = 10_000

integral_approximation = simpson(
    lambda x: math.exp(-x * x),
    -L,
    L,
    n_steps
)

sqrt_pi_g = math.sqrt(pi_g)
difference = integral_approximation - sqrt_pi_g
normalized_area = integral_approximation / sqrt_pi_g

print("Geometric/N-series source values:")
for N, Q in source_values:
    print(f"N = {N:8d}, Q_N = {Q:.17f}")

print()
print("pi_g from geometry/N-series:", format(pi_g, ".17f"))
print("sqrt(pi_g):", format(sqrt_pi_g, ".17f"))
print("Integral approximation:", format(integral_approximation, ".17f"))
print("difference:", difference)
print("normalized area:", format(normalized_area, ".17f"))
\end{lstlisting}

\clearpage

\section{Sample console output}

\begin{lstlisting}
Geometric/N-series source values:
N =    24576, Q_N = 3.14159264503368973
N =    49152, Q_N = 3.14159265145076638
N =    98304, Q_N = 3.14159265305503599

pi_g from geometry/N-series: 3.14159265358979267
sqrt(pi_g): 1.77245385090551588
Integral approximation: 1.77245385090551033
difference: -5.551115123125783e-15
normalized area: 0.99999999999999689
\end{lstlisting}

\begin{thebibliography}{9}

\bibitem{ahlfors}
L. V. Ahlfors,
\emph{Complex Analysis},
McGraw--Hill.

\bibitem{whittakerwatson}
E. T. Whittaker and G. N. Watson,
\emph{A Course of Modern Analysis},
Cambridge University Press.

\bibitem{abramowitzstegun}
M. Abramowitz and I. A. Stegun,
\emph{Handbook of Mathematical Functions},
Dover.

\bibitem{nistdlmf}
NIST Digital Library of Mathematical Functions,
\url{https://dlmf.nist.gov/}.

\end{thebibliography}

\end{document}
